\section{Conclusion}

AI agents are entering high-performance computing, bringing with them a threat category that HPC security was never designed to address: the hijacked authorized agent. Through prompt injection attacks exploiting shared filesystems, co-located compute nodes, compromised agent tools, and coordinated multi-agent exfiltration, an adversary can subvert an agent that already holds valid credentials and legitimate permissions. The hijacked agent exfiltrates data through encrypted, whitelisted LLM API channels, precisely the channels that every legitimate agent must use, making the attack invisible to traditional monitoring and data loss prevention. No existing defense mechanism, including authentication, authorization, intrusion detection, DLP, or user behavior analytics, addresses this threat.

We propose behavioral attestation as the foundation for securing AI agents in HPC. The core insight is pragmatic: we cannot prevent all injection attacks because the attack surface is too large and diverse. However, we can detect and contain their effects by attesting to behavioral constraints. Behavioral attestation provides provable guarantees rather than probabilistic detection, uses constraint-based policies rather than evasion-prone signatures, and enforces containment at runtime rather than post-hoc alerting. Even if an agent is hijacked, the attestation mechanism detects constraint violations and enforces containment before data exfiltration completes.

This paper formalizes four HPC-specific injection attack vectors that have not been studied in prior work, designs and implements AEGIS as a complete zero-trust architecture, and empirically demonstrates its effectiveness. Across all four attack scenarios, AEGIS achieves 100\% detection with sub-second latency, compared to 0-75\% for traditional defenses. The system operates with less than 3\% overhead on AMD EPYC processors, scales to hundreds of agents per node, and produces zero false positives on legitimate HPC workflows. All components are implemented and available as open source.

As AI agents become more autonomous and more deeply integrated into scientific computing infrastructure, the need for provable, runtime behavioral verification will only grow. Behavioral attestation provides the foundation for this verification, a new security primitive for a new class of threats.